% Part: incompleteness
% Chapter: arithmetization-syntax
% Section: coding-terms

\documentclass[../../../include/open-logic-section]{subfiles}

\begin{document}

\olfileid{inc}{art}{trm}
\olsection{Pangodean Suku}

\begin{explain}
Suku iku salah sawijining jinis urutan simbol:
diwangun kanthi induktif saka konstanta lan
variabel miturut aturan pambentukan suku.
Amarga urutan simbol bisa dikodeake minangka
wilangan---nganggo skema pangodean simbol
lan cara ngodeake urutan wilangan---menehi
wilangan G\"odel marang suku ora angel.
Tantangane yaiku nuduhake yen sifat
``sawijining wilangan dadi wilangan G\"odel
saka suku kang kabentuk kanthi bener''
iku komputabel, malah rekursif primitif.
\end{explain}

!!^{variable}s lan !!{constant}s iku suku
kang paling prasaja, lan gampang mriksa apa
$x$ iku wilangan G\"odel saka suku kaya mangkono:
$\fn{Var}(x)$ bener yen $x$ iku $\Gn{\Obj v_i}$
kanggo sawijining~$i$. Tegese, $x$~ngodeake
urutan dawa~$1$ lan unsur tunggale $(x)_0$
iku kode saka !!{variable}~$\Obj v_i$;
yaiku $x$ iku $\tuple{\tuple{1, i}}$
kanggo sawijining~$i$. Semono uga,
$\fn{Const}(x)$ bener yen $x$ iku
$\Gn{\Obj c_i}$ kanggo sawijining~$i$.
Rong relasi iki rekursif primitif,
amarga yen $i$ kaya mangkono ana,
mesthi $< x$:
\begin{align*}
  \fn{Var}(x) & \defiff \bexists{i<x}{x = \tuple{\tuple{1, i}}}\\
  \fn{Const}(x) & \defiff \bexists{i<x}{x = \tuple{\tuple{2, i}}}
\end{align*}

\begin{prop}
\ollabel{prop:term-primrec}
Relasi $\fn{Term}(x)$ lan $\fn{ClTerm}(x)$,
kang bener yen lan mung yen $x$ iku wilangan
G\"odel saka suku utawa suku katutup,
miturut urutane, iku rekursif primitif.
\end{prop}

\begin{proof}
Urutan simbol~$s$ iku suku yen lan mung yen
ana urutan suku~$s_0$, \dots,
$s_{k-1} = s$ kang nyathet carane suku~$s$
diwangun saka !!{constant}s lan !!{variable}s
miturut aturan pambentukan suku. Supaya urutan
kang diduga dadi urutan pambentukan iki
manut aturan, kanggo saben $i < k$ kudu
salah siji saka kahanan iki bener:
\begin{enumerate}
\item $s_i$ iku !!a{variable}~$\Obj v_j$, utawa
\item $s_i$ iku !!a{constant}~$\Obj c_j$, utawa
\item $s_i$ diwangun saka $n$ suku $t_1$,
  \dots, $t_n$ kang wis ana sadurunge
  papan~$i$ nganggo !!{function} $n$ papan
  $\Obj f^n_j$.
\end{enumerate}
Kanggo nuduhake yen relasi kang cocog ing
wilangan G\"odel iku rekursif primitif,
kita kudu nyatakake syarat iki kanthi
rekursif primitif, yaiku nganggo fungsi
lan relasi rekursif primitif uga
kuantifikasi kang winates.

Upamane $y$ iku wilangan kang ngodeake
urutan $s_0$, \dots, $s_{k-1}$,
yaiku $y = \tuple{\Gn{s_0}, \dots, \Gn{s_{k-1}}}$.
Wilangan iku ngodeake urutan pambentukan
kanggo suku kang wilangan G\"odele~$x$
yen lan mung yen kanggo kabeh $i < k$:
\begin{enumerate}
\item $\fn{Var}((y)_i)$, utawa
\item $\fn{Const}((y)_i)$, utawa
\item ana $j$, $n$ lan wilangan~$z = \tuple{z_1, \dots, z_n}$
  saengga saben $z_l$ padha karo sawijining
  $(y)_{i'}$ kanggo $i' < i$, lan
\[
(y)_i = \Gn{\Obj f^n_j(} \concat \fn{flatten}(z) \concat \Gn{)},
\]
\end{enumerate} Cathetan koreksi sumber OLPL-776: klausa katelu saiki ngenalake indeks j kanggo simbol fungsi minangka salah sawijining saksi eksistensial.
lan uga $(y)_{k-1} = x$. (Fungsi
$\fn{flatten}(z)$ ngowahi urutan
$\tuple{\Gn{t_1}, \dots, \Gn{t_n}}$ dadi
$\Gn{t_1, \dots, t_n}$ lan rekursif primitif.)

Indeks $j$, $n$, lan wilangan G\"odel $z_l$ saka suku $t_l$
ing (3) kabeh luwih cilik tinimbang~$y$. Kode~$z$ saka
urutan~$\tuple{z_1, \dots, z_n}$ diwatesi dening
$p_y^{y(y+1)}$. Kita bisa ngganti $k$ ing ndhuwur nganggo
$\len{y}$. Mula pratelan ``$y$ iku kode urutan pambentukan suku
kanthi wilangan G\"odel~$x$'' bisa dites kanthi rekursif primitif.
Cathetan panjelasan SOL6-F072: tes uga mriksa kode urutan sah kang ora kosong.
Wates kode argumentasi aman saka dawa lan saben komponene kang kurang y.

Saiki kita mung perlu mesthekake yen ana
wates rekursif primitif kanggo~$y$.
Yen $x$ iku wilangan G\"odel saka sawijining
suku, mesthi ana urutan pambentukan kanthi
paling akeh $\len{x}$ suku (amarga saben
suku ing urutan pambentukan~$s$ kudu
wiwit ing sawijining papan ing~$s$,
lan ora ana rong subsuku kang wiwit ing
papan kang padha). Wilangan G\"odel saben
subsuku saka~$s$ mesthi $\le x$.
Mula tansah ana urutan pambentukan
kanthi kode $\le p_{k-1}^{k(x+1)}$,
ing ngendi $k=\len{x}$.

Kanggo $\fn{ClTerm}$, cukup tinggalake
klausa kanggo !!{variable}s.
\end{proof}

\begin{prob}
Tuduhna yen fungsi $\fn{flatten}(z)$,
kang ngowahi urutan
$\tuple{\Gn{t_1}, \dots, \Gn{t_n}}$
dadi $\Gn{t_1, \dots, t_n}$,
iku rekursif primitif.
\end{prob}

\begin{prop}\ollabel{prop:num-primrec}
  Fungsi $\fn{num}(n) = \Gn{\num{n}}$
  iku rekursif primitif.
\end{prop}

\begin{proof}
  Kita netepake $\fn{num}(n)$
  kanthi rekursi primitif:
  \begin{align*}
    \fn{num}(0) & = \Gn{\Obj 0}\\
    \fn{num}(n+1) & = \Gn{\prime(} \concat \fn{num}(n) \concat \Gn{)}.
  \end{align*}
\end{proof}

\end{document}
